% SOURCE_AUTHORITY: sga5_fr_workpass.tex, lines 4796--4833 (LF-normalized unit).
% SOURCE_SHA256: 791F4EFFC5E02832D5D77ED03518C8156D6F07E4C8238B03545DB93D883FBB28
La demostración se apoyará en el

\begin{statement}{Lema 4.3}
Sean $Z$ un esquema noetheriano regular conexo de dimensión $2$, de características residuales coprimas con $\ell$, y $X$, $Y$ divisores regulares que se cortan transversalmente en un punto $z$, de donde se obtiene un diagrama de inclusiones
\[
\begin{tikzcd}[column sep=large,row sep=large]
Z-(X\cup Y) \arrow[r,hook,"i'"] \arrow[d,hook,"j'"'] & Z-Y \arrow[d,hook,"j"]\\
Z-X \arrow[r,hook,"i"'] & Z\\
Y-z \arrow[r,hook,"i_1"'] \arrow[u,hook,"m'"'] & Y \arrow[u,hook,"m"] .
\end{tikzcd}
\]
Se adopta la hipótesis de pureza:
\[
(P)\qquad \R^qj_*(\mathbb Z/\ell)=0\quad\text{para }p>2.
\]
Sea, por otra parte, $L\in\Ob D^b(Z-(X\cup Y),\Lambda)$, con cohomología constructible y localmente constante, cuya representación correspondiente de $\pi_1(Z-(X\cup Y))$ se factorice a través de un $\ell$-grupo. Entonces:
\begin{enumerate}[label=\alph*)]
\item Se tiene un isomorfismo canónico
\[
i_{1!}m'^*\R j'_*L\xrightarrow{\sim} m^*\R j_*(i'_!L),
\]
y el complejo $m'^*\R j'_*L$ tiene cohomología constructible y localmente constante.
\item Supongamos que $Y$ es conexo, y sea $\bar s$ un punto geométrico de $Y$ situado sobre un punto cerrado $s\ne z$. Entonces la flecha de restricción
\begin{equation}
\R\Gamma(Y,m^*\R j_*(i'_!L))\longrightarrow (\R j_*(i'_!L))_{\bar s}
\tag{4.3.1}
\end{equation}
es nula.
\item Bajo las hipótesis de b), sea $C$ un subesquema cerrado de dimensión $1$ de $Z$ tal que $Y\cap C=s$, y denotemos por $\widetilde C$ la localización estricta de $C$ en $\bar s$. Entonces la flecha de restricción
\begin{equation}
\R\Gamma(Z-Y,i'_!L)
\longrightarrow
\R\Gamma(\widetilde C-\bar s,i'_!(L)|\widetilde C-\bar s)
\tag{4.3.2}
\end{equation}
es nula.
\end{enumerate}
\end{statement}
